\scanpage{5}
[On est \unsure{amené|ramené} à prouver ceci : Soit $\psi:Y\to M$ une application telle que (i) pour \deleted{tt} $y\in Y$, $k(\psi(y))=k(y)$, i.e. $\psi(y)$ est rationnel sur $k(y)$ (ii) l'application correspondante $\varphi:Y'\to X'$, \deleted{\unsure{est une}} faisant correspondre à tout $y\in Y'$ la composante \unsure{ult.|rel\textsuperscript{t}} irréd. de $f^{-1}(y)$ correspondant à \deleted{$k(y)$} $\psi(y)$, est \unsure{monisante|monisente|monogène}. -- Alors, \deleted{\unsure{que c'est une}} $\psi(Y)$ est \deleted{à la fois} ouvert. \deleted{et fermé, \unsure{et} $\psi(Y)\to Y$ est} donc $\psi$ provient d'une section de $M$ sur $Y$]~; sous cette forme, c'est évident, grâce à la condition de continuité mise \unsure{dans|ds} la définition de ``monisant'']

\uline{Conséquence} Si le Pb des modules a une solution pour $X\times_Y\operatorname{Spec}(\mathcal O_y)$ \inserted{sur $\operatorname{Spec}(\mathcal O_y)$}, \deleted{\unsure{elle a une}} il existe un ouvert $U\ni y$ tel que le Pb des modules ait une solution pour $X|U\,/\,U$.

Cela nous ramène, pour l'existence, au cas où \deleted{$X$} $Y=\operatorname{Spec}(A)$, $A$ local. \inserted{[Car \deleted{\unsure{tt}} alors supposons \unsure{$\dim Y=n<+\infty$}]} D'ailleurs, grâce à la descente des morphismes étales, on est ramené au cas où $A$ est \uline{complet},
% slanted
\margin{? descente des morphismes étales dans le cas non séparé \deleted{\unsure{finis}} ???}
et on peut $=$ supposer que la fibre $f^{-1}(y_0)$ de son point fermé se décompose en composantes \uline{ult\textsuperscript{t}} irréductibles.
% archive sheet number: 4
